% Auto-typeset from the previously authored formula-by-formula study manuscript.
% Based on the thesis main branch f80e5ac903dd, 2026-10-31.

\chapter{Jost 解、解析性、散射矩阵与共振相容性}

对照正文：../hmb-nzbc-thesis/section/03-scattering.tex  \par
覆盖 §3.1–§3.4，特别是 thm:jost、prop:scattering、prop:quartet、prop:evolution、prop:resonance、prop:highenergy。\par

\textbf{学习目标}：能从二阶 AKNS 系统独立写出 Volterra 方程，判断四个 Jost 列的解析域，使用 Cramer 法则求出散射系数，并说明共振点反射为零的具体前提。\par

\section{第一段：为什么先固定 z？}

HMB 的 x 方程用于定义散射数据，z 方程用于决定散射数据如何传播。最初构造 Jost 解时，我们可以把 z 固定为一个参数，研究\par

\[
\chi_x=\widehat U(x,z,\lambda)\chi,
\qquad
\widehat U=-\mathrm{i}\lambda\sigma_3+A(u).
\]

当 $x\to{}\pm{}\infty{}$ 时 $u\to{}Q$，故 $\widehat U\to{}U_{\infty}$。把两者之差记为\par

\[
\Delta A(x)=A(u(x))-A(Q)
=\begin{pmatrix}0&u-Q\\-(\overline u-Q)&0\end{pmatrix}.
\]

若 $u-Q\in{}L^{1}\cap{}L\infty{}$，这个扰动沿实 x 可积。$L^{1}$ 保证 Volterra 积分的总强度有限；$L\infty{}$ 帮助控制局部乘积和矩阵核。\par

注意这时\textbf{不是假设 q 本身在无穷远趋零}，而是减掉常数 Q 后的扰动可积。这是非零背景与第五章零背景的根本差别。\par

\section{Jost 归一化为什么要乘一个指数？}

背景常系数系统的基本解为 $E(w)e^{-i\rho{}x\sigma{}_{3}}$。所以自然定义从左、右逼近背景的空间基本解：\par

\[
\chi_\pm(x,w)e^{\mathrm{i}\rho x\sigma_3}
\longrightarrow E(w),\qquad x\to\pm\infty.
\]

把指数移到矩阵内部定义\par

\[
\mu_\pm=\chi_\pm e^{\mathrm{i}\rho x\sigma_3}.
\]

这样 $\mu_{\pm}$ 的端态是 E，而不是快速振荡的 $E e^{-i\rho{}x\sigma{}_{3}}$；后续把左右两套解析列拼成 RHP 更方便。\par

下标 $\pm{}$ 表示归一化端点，不是取复平面的上下边界。$\mu{}_{+,1}$ 是右 Jost 矩阵第一列，$\mu{}_{-,2}$ 是左 Jost 矩阵第二列。\par

\section{thm:jost 和 eq:volterra：完整的变分常数推导}

定义 $W_{\pm}=E^{-1}\chi_{\pm}$。利用 $U_{\infty}E=-i\rho{}E\sigma{}_{3}$，原方程给出\par

\[
\begin{aligned}
(W_\pm)_x
&=E^{-1}\chi_{\pm,x}\\
&=E^{-1}(U_\infty+\Delta A)E W_\pm\\
&=-\mathrm{i}\rho\sigma_3 W_\pm
+E^{-1}\Delta A E W_\pm.
\end{aligned}
\]

先在有限起点 $x_{0}$ 应用变分常数：\par

\[
W_\pm(x)
=e^{-\mathrm{i}\rho(x-x_0)\sigma_3}W_\pm(x_0)
+\int_{x_0}^{x}e^{-\mathrm{i}\rho(x-y)\sigma_3}
E^{-1}\Delta A(y)E W_\pm(y)\,\mathrm{d}y.
\]

左右乘回 E 和 $e^{i\rho{}x\sigma{}_{3}}$，再分别令 $x_{0}\to{}-\infty{}$ 或 $+\infty{}$。最终得原文 eq:volterra：\par

\[
\boxed{
\begin{aligned}
\mu_\pm(x,w)=E(w)
+\int_{\pm\infty}^{x}
&E(w)e^{-\mathrm{i}\rho(x-y)\sigma_3}E(w)^{-1}\\
&\cdot\Delta A(y)\mu_\pm(y,w)
e^{\mathrm{i}\rho(x-y)\sigma_3}\,\mathrm{d}y.
\end{aligned}}
\]

为什么右侧最后还有一个右乘指数？因为 $\mu{}=\chi{} e^{i\rho{}x\sigma{}_{3}}$；右侧乘指数不可以与左侧矩阵自由交换。$e^{i\rho{}x\sigma{}_{3}}$ 是对角矩阵，左、右两次指数属于不同乘法位置。\par

“从 $+\infty{}$ 到 x 的积分”按\par

\[
\int_{+\infty}^{x}f(y)\,\mathrm{d}y
=-\int_x^{+\infty}f(y)\,\mathrm{d}y
\]

解释。这决定后面哪一个指数在积分区域内衰减。\par

\section{逐列计算指数有界性}
\noindent{\small\textbf{对应原文：}\texttt{eq:columnkernel}}\par

先设\par

\[
\sigma_3=\operatorname{diag}(1,-1),\qquad
s_1=1,\quad s_2=-1.
\]

观察 $\mu{}$ 的第 j 列：右侧 $e^{i\rho{}(x-y)\sigma{}_{3}}$ 作用在 $e_j$ 上，只给它乘一个标量 $e^{i s_j \rho{}(x-y)}$。于是核简化为\par

\[
K_j(x-y,w)=e^{\mathrm{i}s_j\rho(x-y)}
E e^{-\mathrm{i}\rho(x-y)\sigma_3}E^{-1}.
\]

对 $j=1$：\par

\[
e^{\mathrm{i}\rho(x-y)}
e^{-\mathrm{i}\rho(x-y)\sigma_3}
=\operatorname{diag}(1,e^{2\mathrm{i}\rho(x-y)}).
\]

对 $j=2$：\par

\[
e^{-\mathrm{i}\rho(x-y)}
e^{-\mathrm{i}\rho(x-y)\sigma_3}
=\operatorname{diag}(e^{-2\mathrm{i}\rho(x-y)},1).
\]

假设 $w\in{}D_{+}$，即 $Im\rho{}>0$。\par

\begin{itemize}[leftmargin=2em]
\item 若选择负端归一化 \ensuremath{\chi}\textsubscript{-}，积分变量 y\ensuremath{\le}x，故 x\ensuremath{-}y\ensuremath{\ge}0；对第一列 |e\textasciicircum{}\{2i\ensuremath{\rho}(x\ensuremath{-}y)\}|\ensuremath{\le}1。\textbf{\ensuremath{\mu}\_\{-,1\} 在 D\textsubscript{+} 解析}。
\item 若选择正端归一化 \ensuremath{\chi}\textsubscript{+}，积分变量 y\ensuremath{\ge}x，故 x\ensuremath{-}y\ensuremath{\le}0；对第二列 |e\textasciicircum{}\{\ensuremath{-}2i\ensuremath{\rho}(x\ensuremath{-}y)\}|\ensuremath{\le}1。\textbf{\ensuremath{\mu}\_\{+,2\} 在 D\textsubscript{+} 解析}。
\end{itemize}

如果 $Im\rho{}<0$，上述指数不再按同一选择有界，需要交换列与端点。结论为\par

\[
\begin{array}{c|c}
D_+:\operatorname{Im}\rho>0 & \mu_{-,1},\ \mu_{+,2}\\
D_-:\operatorname{Im}\rho<0 & \mu_{+,1},\ \mu_{-,2}
\end{array}
\]

这里是帮助理解的二维表，不表示完整矩阵在整个 $D_{+}$ 或 $D_{-}$ 同时解析。\par

\section{为何 Neumann 级数收敛？1/n! 从何而来？}

以负端第一列为例，设积分算子为 K，则\par

\[
\mu_{-,1}=E_1+KE_1+K^2E_1+\cdots.
\]

第 n 次迭代的积分区域不是整个 n 维立方体，而是有序单纯形\par

\[
-\infty<y_n<\cdots<y_2<y_1<x.
\]

若核有界于 $C_K\parallel{}\Delta{}A(y)\parallel{}$，由积分对称性得到\par

\[
\begin{aligned}
\|(K^nE_1)(x)\|
&\le\|E_1\|C_K^n
\int\limits_{y_n<\cdots<y_1<x}
\prod\limits_{\ell=1}^{n}\|\Delta A(y_\ell)\|
\,\mathrm{d}y_1\cdots\mathrm{d}y_n\\
&=\dfrac{\|E_1\|C_K^n}{n!}
\left(\int_{-\infty}^{x}\|\Delta A(y)\|\,\mathrm{d}y\right)^n.
\end{aligned}
\]

最后一个等号因为对称被积函数在 n! 个变量排列区域上的积分相等。有阶乘分母，级数绝对收敛；同样估计作用于两解之差，说明唯一性。每次迭代项都随 w 解析，紧集一致收敛由 Weierstrass 定理保证解析性。\par

\textbf{关键限制}：$E^{-1}$ 含 $1/\gamma{}(w)$，在 $w=\pm{}iQ$ 退化，不能直接把此证明延伸到分支点。分支点可能需要带权的空间衰减假设与单独的阈值展开。\par

\section{散射矩阵究竟表示什么？}
\noindent{\small\textbf{对应原文：}\texttt{eq:Sdef}}\par

$\chi{}_{+}$、$\chi{}_{-}$ 是同一二阶空间系统的两套可逆基本矩阵，所以\par

\[
S_{\mathrm{sp}}(w,z)=\chi_+(x,w)^{-1}\chi_-(x,w)
\]

与 x 无关。原文约定\par

\[
\chi_-=\chi_+S_{\mathrm{sp}}.
\]

这是一个重要的方向约定。如果改成 $\chi{}_{+}=\chi{}_{-}S$，所有反射系数及 RHP 跳跃符号都要相应改变。\par

为了使用完整传播相位，定义同时归一化的基本矩阵\par

\[
\Psi_\pm=\chi_\pm e^{\mathrm{i}\rho Fz\sigma_3},
\qquad
\Psi_-=\Psi_+S.
\]

S 与 $S_sp$ 的区别：S 是两部分 Lax 对同时演化的散射矩阵，在后面证明的正则性条件下与 z 无关；$S_sp$ 则可随 z 通过对角相位共轭变化。\par

\section{为什么散射矩阵行列式为 1？}
\noindent{\small\textbf{对应原文：}\texttt{eq:detjost}}\par

对基本矩阵 $\chi{}_{x}=U\chi{}$，有 Liouville（Jacobi）行列式公式\par

\[
\dfrac{\mathrm{d}}{\mathrm{d}x}\det\chi
=\operatorname{tr}U\det\chi.
\]

这里 $tr U=0$，所以 $det\chi{}$ 与 x 无关；利用两端归一化\par

\[
\det\chi_\pm=\det E=\gamma(w).
\]

于是\par

\[
\det S_{\mathrm{sp}}
=\dfrac{\det\chi_-}{\det\chi_+}=1.
\]

同时归一化右乘的对角指数行列式也为 1，所以 $detS=1$。这个等式是以后计算跳跃矩阵 $detJ=1$ 的原因之一。\par

\section{Cramer 法则逐列推导}
\noindent{\small\textbf{对应原文：}\texttt{eq:wronskian}}\par

把 $\Psi{}_{-}=\Psi{}_{+}S$ 写成列：\par

\[
\begin{aligned}
\Psi_{-,1}&=s_{11}\Psi_{+,1}+s_{21}\Psi_{+,2},\\
\Psi_{-,2}&=s_{12}\Psi_{+,1}+s_{22}\Psi_{+,2}.
\end{aligned}
\]

对第一式两边与 $\Psi{}_{+},_{2}$ 构成行列式：\par

\[
\det(\Psi_{-,1},\Psi_{+,2})
=s_{11}\det(\Psi_{+,1},\Psi_{+,2}),
\]

因为包含两列 $\Psi{}_{+},_{2}$ 的那项为零。分母 $det\Psi{}_{+}=\gamma{}$，得到\par

\[
s_{11}=\dfrac{\det(\Psi_{-,1},\Psi_{+,2})}{\gamma}.
\]

同理，与 $\Psi{}_{+},_{1}$ 构成适当的行列式，得到\par

\[
\begin{aligned}
s_{22}&=\dfrac{\det(\Psi_{+,1},\Psi_{-,2})}{\gamma},\\
s_{21}&=\dfrac{\det(\Psi_{+,1},\Psi_{-,1})}{\gamma},\\
s_{12}&=\dfrac{\det(\Psi_{-,2},\Psi_{+,2})}{\gamma}.
\end{aligned}
\]

最后一个分子列序不能交换，否则会多一个负号。\par

由于 $s_{11}$ 的分子是 $D_{+}$ 的两列解析函数的行列式，且 $\gamma{}$ 在普通点非零，$s_{11}$ 可以延拓进 $D_{+}$；$s_{22}$ 类似在 $D_{-}$ 解析。反射相关的 $s_{12}$、$s_{21}$ 一般只在连续谱定义，不应无条件宣称它们在某个开域解析。\par

\section{eq:conjugation、eq:sconj：为何存在复共轭对称？}

定义\par

\[
K=\begin{pmatrix}0&1\\-1&0\end{pmatrix}.
\]

空间 Lax 矩阵的实参数约化意味着，如果 $\Psi{}(x,z,\overline w)$ 是某个基本解，把它逐项共轭后再做 K 共轭，也产生具有相应归一化的基本解。Jost 解唯一，故\par

\[
\Psi_\pm(w)=K\overline{\Psi_\pm(\overline w)}K^{-1}.
\]

对散射矩阵也有相同形式。若 $S=\begin{pmatrix}s_{11}&s_{12}\\s_{21}&s_{22}\end{pmatrix}$，直接相乘可得到\par

\[
s_{22}(w)=\overline{s_{11}(\overline w)},
\qquad
s_{12}(w)=-\overline{s_{21}(\overline w)}.
\]

这一对称把 $D_{+}$ 的透射零点送到 $D_{-}$ 的共轭零点。它与换叶对称是两种不同操作，不可以混为一谈。\par

\section{eq:Esheet、eq:musheet：换叶与共轭不是一回事}

从第二章\par

\[
h(w)=-\dfrac{Q^2}{w},\quad
\lambda(h(w))=\lambda(w),\quad \rho(h(w))=-\rho(w)
\]

出发，对 E 直接计算可得\par

\[
E(h(w))=E(w)C(w),
\qquad
C(w)=\dfrac{\mathrm{i}w}{Q}\sigma_1.
\]

$C \sigma{}_{3}=-\sigma{}_{3} C$，正好把两种背景指数模式交换。于是\par

\[
\Psi_\pm(h(w))=\Psi_\pm(w)C(w),
\]

并由 $\Psi{}_{-}=\Psi{}_{+}S$ 导出\par

\[
S(h(w))=\sigma_1S(w)\sigma_1.
\]

后者相当于交换透射和反射矩阵的列、行位置。请动手算 $\sigma{}_{1}S\sigma{}_{1}$，确认它确实使 $s_{11}\leftrightarrow{}s_{22}$、$s_{12}\leftrightarrow{}s_{21}$。\par

\textbf{与共轭区别}：h 保持 $\lambda{}$ 不变但改变 $\rho{}$ 符号；复共轭通常把 $\lambda{}$、$\rho{}$ 变成共轭值，改变相应解析域。两者共同生成四点轨道。\par

\section{反射系数只是比例}
\noindent{\small\textbf{对应原文：}\texttt{eq:reflection}}\par

原文采用\par

\[
r=\dfrac{s_{21}}{s_{11}},
\qquad
\widetilde r=\dfrac{s_{12}}{s_{22}}.
\]

这是散射矩阵两个非对角系数相对于对应透射系数的比值。若 $s_{11}$ 或 $s_{22}$ 在连续谱为零，比例就可能没有定义，因此论文额外假定普通连续谱无透射零点。\par

反射为零时 $r=\widetilde r=0$，意味着 S 在普通连续谱是对角矩阵，并不意味着 $S=I$ 或没有离散零点。第五章的 S 恰为一个非平凡的对角矩阵。\par

\section{prop:quartet 与 eq:quartet：四点离散谱的来历}

若 $\zeta{}\in{}D_{+}$ 且 $s_{11}(\zeta{})=0$，那么由共轭对称得到 $s_{22}(\overline \zeta{})=0$；由换叶对称得到 $s_{22}(h(\zeta{}))=0$；再组合两种对称得到 $s_{11}(h(\overline \zeta{}))=0$。四点为\par

\[
\left\{\zeta,\overline\zeta,-\dfrac{Q^2}{\zeta},-\dfrac{Q^2}{\overline\zeta}\right\}.
\]

其中 $\zeta{}$、$h(\overline \zeta{})$ 在 $D_{+}$，另外两个在 $D_{-}$。一般简单零点的四点互不相同。第五章外部纯虚谱点 iW 的 $D_{+}$ 配对点是 $-iQ^{2}/W$。\par

为什么透射零点对应本征函数？在 $s_{11}(\zeta{})=0$ 处，\par

\[
\det(\Psi_{-,1}(\zeta),\Psi_{+,2}(\zeta))=0.
\]

两列线性相关，可以形成同一向量解：它向左、右分别由所选的衰减模态控制，因而与局域本征态相联系。线性相关常数 $b_j$ 是“范数常数”的原始来源。\par

\section{\texorpdfstring{$b_j$、$c_j$ 与 $C_j$}{bj, cj 与 Cj} 各表示什么？}

原文定义\par

\[
\Psi_{-,1}(\zeta_j)=b_j\Psi_{+,2}(\zeta_j),
\qquad
c_j=\dfrac{b_j}{s_{11}'(\zeta_j)}.
\]

由于 $s_{11}$ 在 $\zeta{}_j$ 有简单零点，$\dfrac{1}{s}_{11}$ 的留数为 $\dfrac{1}{s}'_{11}(\zeta{}_j)$。当第一列被除以 $s_{11}$ 后，留数系数就出现 $c_j$。去振荡后的列还引入\par

\[
C_j(x,z)=c_j e^{2\mathrm{i}\theta(\zeta_j;x,z)}.
\]

请区分：$c_j$ 是在正确的同时归一化约定下与 x、z 无关的常数；$C_j$ 由于含 $\theta{}$，仍然随 x、z 变化。在第五章有限维系统里真正出现的是 $C_j$，不是把 $c_j$ 直接代进去。\par

\section{两种散射矩阵为什么演化方式不同？}
\noindent{\small\textbf{对应原文：}\texttt{prop:evolution}}\par

背景时间矩阵为 $V_{\infty}=-FU_{\infty}$，对角模式传播含 $e^{i\rho{}Fz\sigma{}_{3}}$。定义 $\Psi_{\pm}=\chi_{\pm}e^{i\rho{}Fz\sigma{}_{3}}$ 之后，若 $\Psi_{\pm}$ 同时满足两部分 Lax 方程，比较同一组基本解得到 S 不依赖 x、z。\par

由 $\chi{}_{-}=\chi{}_{+}S_sp$ 与 $\Psi{}_{-}=\Psi{}_{+}S$，代入 $\chi{}=\Psi{}e^{-i\rho{}Fz\sigma{}_{3}}$，整理：\par

\[
S_{\mathrm{sp}}(w,z)
=e^{\mathrm{i}\rho Fz\sigma_3}S(w)
e^{-\mathrm{i}\rho Fz\sigma_3}.
\]

写 $E_z=\operatorname{diag}(e^{i\rho{}Fz},e^{-i\rho{}Fz})$，则\par

\[
E_z\begin{pmatrix}s_{11}&s_{12}\\s_{21}&s_{22}\end{pmatrix}E_z^{-1}
=\begin{pmatrix}
s_{11}&e^{2\mathrm{i}\rho Fz}s_{12}\\
e^{-2\mathrm{i}\rho Fz}s_{21}&s_{22}
\end{pmatrix}.
\]

因此对角透射系数不变，而空间归一化反射系数满足原文 eq:timer：\par

\[
r_{\mathrm{sp}}(w,z)
=e^{-2\mathrm{i}\rho Fz}r_{\mathrm{sp}}(w,0),
\quad
\widetilde r_{\mathrm{sp}}(w,z)
=e^{2\mathrm{i}\rho Fz}\widetilde r_{\mathrm{sp}}(w,0).
\]

对 S 不依赖 z 的证明，用零曲率对 $\chi{}_{\pm{},z}-V\chi_{\pm}$ 求 x 导数，识别其端态，再利用 Jost 唯一性。这一步要求对 z 微分合法，因此论文额外引入 $W^{3,1}$、$C^{1}$ 等充分正则性。不能把这些演化公式无条件用于所有粗糙数据。\par

\section{为什么在共振点 \texorpdfstring{$r_{\mathrm{sp}}(a)=0$}{rsp(a)=0}？}
\noindent{\small\textbf{对应原文：}\texttt{prop:resonance}}\par

这不是从“F 在 a 发散”直接得出的，而是从介质方程严格推导：\par

\[
R_x=[\widehat U(a),R].
\]

对共振处可逆的空间基本解 $\chi{}$，利用 $\chi{}_{x}=\widehat U(a)\chi{}$：\par

\[
\begin{aligned}
(\chi^{-1}R\chi)_x
&=-\chi^{-1}\widehat U(a)R\chi
+\chi^{-1}R_x\chi
+\chi^{-1}R\widehat U(a)\chi\\
&=\chi^{-1}\bigl(R_x-[\widehat U(a),R]\bigr)\chi=0.
\end{aligned}
\]

所以 $\chi{}^{-1}R\chi{}$ 沿 x 为常数。左右端态及 E(a) 的对角化，确定该常矩阵都是\par

\[
d_a\sigma_3,\qquad
d_a=\dfrac{\eta_0\rho(w_a)}{a}.
\]

利用 $\chi{}_{-}=\chi{}_{+}S_sp$：\par

\[
S_{\mathrm{sp}}^{-1}d_a\sigma_3S_{\mathrm{sp}}
=d_a\sigma_3.
\]

因为 $d_{a}\ne{}0$，等价于 $[\sigma{}_{3},S_sp]=0$。设 $S_sp=\begin{pmatrix}a&b\\c&d\end{pmatrix}$，则\par

\[
[\sigma_3,S_{\mathrm{sp}}]
=\begin{pmatrix}0&2b\\-2c&0\end{pmatrix}=0
\ \Longrightarrow\ b=c=0.
\]

所以反射系数在共振点为零。\par

\textbf{所有前提都必须记住}：固定 z，介质方程成立；两端有相容背景；Jost 基在 a 正则可逆；$\eta{}_{0}\ne{}0$；透射分母非零。这不是一般无条件结论，也不是“$r(a)=0$ 足以保证 RHP 一般可解”。\par

对于不同介质端态，左右常矩阵可能不同，得到的是相似关系而非与 $\sigma{}_{3}$ 对易。\par

\section{为什么需要做高能展开？}
\noindent{\small\textbf{对应原文：}\texttt{prop:highenergy}}\par

第四章需要 $M(w)=I+M^{(1)}/w+O(w^{-2})$，第五章要从 $M^{(1)}$ 读取光场 u。光说“由渐近可得”不足以计算系数，因此本章在额外条件 $u-Q\in{}W^{2,1}$ 下推导高能展开。\par

原文引入与空间势对应的 Hermitian 矩阵\par

\[
B(u)=\begin{pmatrix}0&u\\\overline u&0\end{pmatrix},
\qquad A(u)=\sigma_3B(u).
\]

及累计积分\par

\[
J_\pm(x)=\int_{\pm\infty}^{x}(Q^2-|u(y)|^2)\,\mathrm{d}y.
\]

为什么这个积分是有限的？因为 u-Q 可积而 u 有界，且\par

\[
Q^2-|u|^2
=(Q-u)\overline u+Q(\overline Q-\overline u)
\]

在 Q 实值时可由 |u-Q| 控制，因此可积。注意 $J_\pm{}$ 端点不同，后面差值变成全实轴积分。\par

论文高能结论 eq:mu-highenergy：\par

\[
\mu_\pm=I+\dfrac{-\mathrm{i}B(u)+\mathrm{i}J_\pm\sigma_3}{w}
+O(w^{-2}),
\qquad w\to\infty.
\]

这里公式矩阵式需按具有相应解析性的列理解。\par

\section{为什么要设 \texorpdfstring{$T=I-\dfrac{\ii B}{w}$}{T=I-iB/w}？}

直接从原积分方程展开，背景 E 与非对角指数同时出现，第一阶系数不易提取。做“局部规范”：\par

\[
T(x,w)=I-\dfrac{\mathrm{i}B(u(x))}{w},
\qquad \mu_\pm=Tn_\pm.
\]

由于 $B^{2}=|u|^{2}I$，\par

\[
\begin{aligned}
T\left(I+\dfrac{\mathrm{i}B}{w}\right)
&=I+\dfrac{B^2}{w^2}
=\left(1+\dfrac{|u|^2}{w^2}\right)I,\\
T^{-1}
&=\dfrac{I+\mathrm{i}B/w}{1+|u|^2/w^2}.
\end{aligned}
\]

原文 n-volterra-system 最关键的一步是将 $UT+i\rho{}T\sigma{}_{3}-T_{x}$ 化简。使用 $A=\sigma{}_{3}B$、$B\sigma{}_{3}=-\sigma{}_{3}B$、$\lambda{}+\rho{}=w$ 和 $\rho{}-\lambda{}=Q^{2}/w$，得到\par

\[
UT+\mathrm{i}\rho T\sigma_3-T_x
=\dfrac{\mathrm{i}}{w}\bigl((Q^2-|u|^2)\sigma_3+B_x\bigr).
\]

所以 n 满足\par

\[
(n_\pm)_x=-\mathrm{i}\rho[\sigma_3,n_\pm]
+\mathcal Cn_\pm,
\]

其中\par

\[
\mathcal C=\dfrac{1}{w}T^{-1}
\,\mathrm{i}\bigl((Q^2-|u|^2)\sigma_3+B_x\bigr).
\]

由于可积的扰动乘了 1/w，积分迭代自然得到 $n-I=O(\dfrac{1}{w})$。\par

\section{第一迭代如何抽取 \texorpdfstring{$J_\pm$}{J+-}，非对角项为什么更小？}

第一次迭代中对角项没有快速振荡，主项就是\par

\[
\dfrac{\mathrm{i}}{w}\int_{\pm\infty}^{x}
(Q^2-|u(y)|^2)\,\mathrm{d}y
=\dfrac{\mathrm{i}J_\pm(x)}{w}.
\]

非对角项会出现\par

\[
\dfrac{1}{w}\int_{-\infty}^{x}
e^{2\mathrm{i}\rho(x-y)}\overline{u_y(y)}\,\mathrm{d}y.
\]

以积分内核做分部积分（不引入机械 u、v 记号）：\par

\[
\begin{aligned}
&\int_{-\infty}^{x}
e^{2\mathrm{i}\rho(x-y)}\overline{u_y(y)}\,\mathrm{d}y\\
&=-\dfrac{\overline{u_x(x)}}{2\mathrm{i}\rho}
+\dfrac{1}{2\mathrm{i}\rho}
\int_{-\infty}^{x}
e^{2\mathrm{i}\rho(x-y)}\overline{u_{yy}(y)}\,\mathrm{d}y.
\end{aligned}
\]

因为在允许的解析区域指数模$\le{}1$、$u_{x}$ 有界、$u_{xx}$ 可积，右侧为 $O(1/\rho{})=O(\dfrac{1}{w})$。外面再乘 1/w，得到 $O(w^{-2})$；这就解释第一阶非对角修正实际上已由 T 中的 -iB/w 捕获。\par

\section{透射高能系数和第二个无穷远}

由 $s_{11}=det(\mu{}_{-},_{1},\mu{}_{+},_{2})/\gamma{}$，把两列展开，得到\par

\[
s_{11}(w)=1+\dfrac{\mathrm{i}J_{\mathrm{tot}}}{w}
+O(w^{-2}),
\]

其中\par

\[
J_{\mathrm{tot}}=
\int_{-\infty}^{+\infty}(Q^2-|u|^2)\,\mathrm{d}x.
\]

理由是 $J_{-}(x)-J_{+}(x)$ 等于整个实轴积分，与 x 无关。由于 $\gamma{}=1+O(w^{-2})$，不会改变 1/w 的系数。\par

换叶 $h(w)=-Q^{2}/w$ 将大 w 送到小 w。把大 w 的 Jost 展开和换叶对称结合，可得 eq:smallmu：\par

\[
\mu_\pm(w)=-\dfrac{\mathrm{i}Q}{w}\sigma_1+O(1),
\qquad
S(w)=I+O(w),\qquad w\to0.
\]

这说明非零背景 RHP 有\textbf{两种归一化}：$w\to{}\infty{}$ 时趋于 I，$w\to{}0$ 时有指定 1/w 主项。第四章的部分分式式因此必须包含固定的 $-iQ\sigma{}_{1}/w$ 项。\par

\section{自测与阅读关卡}

\begin{enumerate}[leftmargin=2em]
\item 不看正文，自己从 W=E\textsuperscript{-}\textsuperscript{1}\ensuremath{\chi} 推导 Volterra 方程。
\item 判断 e\textasciicircum{}\{2i\ensuremath{\rho}(x\ensuremath{-}y)\} 在 x\ensuremath{-}y\ensuremath{\ge}0、Im\ensuremath{\rho}>0 下的模。
\item 用 Cramer 法则独立求出 s\textsubscript{1}\textsubscript{1} 和 s\textsubscript{1}\textsubscript{2}，并检查列顺序的负号。
\item 把 S\_sp 的对角相位共轭做完整 2\ensuremath{\times}2 乘法，得出两个反射系数的指数。
\item 用两次对称操作找出 \ensuremath{\zeta} 的四点轨道。
\item 写出 r\_sp(a)=0 定理的全部条件，特别是左右介质端态和 Jost 正则性。
\item 亲自完成 T\textsuperscript{-}\textsuperscript{1} 的运算与一次分部积分，说明 O(w\textsuperscript{-}\textsuperscript{2}) 从哪里来。
\end{enumerate}

读懂这一章后，第四章的 RHP 矩阵不再是某种“凭空定义”的未知函数，而是我们已经知道解析域的 Jost 列的自然拼接。\par

